\begin{eg}[AST and Proof Rules]
Consider Language ADDITION. We add a meta-variable \(a\) that ranges over strings generated by the grammar \(S \Coloneqq z \mid S' + S'\): 
\[
  a \in S' \Coloneqq z \mid a_1 + a_2 \quad\quad \text{where } + \text{ is right-associative} 
\]
We then try to define an ``evaluation step'' with a proof system whose statements (judgments) look like
\[
  a_1 \to a_2
\]
We have the inference rules:
\[
  \textsc{(PLUS)}\ z_1 + z_2 \to z_3 \quad\text{if}\quad z_3 = z_1 + z_2
\]
\[
  \textsc{(LEFT)}\ \dfrac{a_1 \to a_1^{\prime}}{a_1 + a_2 \to a_1^{\prime} + a_2}
\]
\[
  \textsc{(RIGHT)}\ \dfrac{a \to a^{\prime}}{z + a \to z + a^{\prime}}
\]
The program will then evaluate to an integer \(z\). Now consider evaluating \(1 + 2 + 3\). 

For \textsc{PLUS}, we have \(z_1 + z_2\); however, we cannot instantiate meta-variables in \(z_1 + z_2\) to get \(1 + 2 + 3\), since both \(z_1\) and \(z_2\) must be integers (single-node ASTs). 

However, for \textsc{LEFT} and \textsc{RIGHT}, we can match \(1 + 2 + 3\). For \textsc{LEFT}, we have \(a_1 + a_2\), where \(a_i\) ranges over any expression in the language, so any sub-AST can replace \(a_i\). This is the same for \textsc{RIGHT}. 

Consider \textsc{LEFT}. After instantiating \(a_1 = 1\) and \(a_2 = 2 + 3\), we have 
\[
  \dfrac{1 \to a_1^{\prime}}{1 + 2 + 3 \to a_1^{\prime} + 2 + 3}
\]
To finish this proof, we need to prove the premise \(1 \to a_1^{\prime}\) for some instantiation of \(a_1^{\prime}\). However, no rule applies to \(1\), since it is already a value, so the proof fails. 

Next, consider \textsc{RIGHT}:
\[
  \dfrac{2 + 3 \to a^{\prime}}{1 + 2 + 3 \to 1 + a^{\prime}}
\]
To finish the proof, we need to prove \(2 + 3 \to a^{\prime}\). This does not fail; instead, we must continue evaluating \(2 + 3\). 

Now we can apply \textsc{PLUS} by instantiating \(z_1 = 2\) and \(z_2 = 3\):
\[
  2 + 3 \to z_3 \quad\text{if}\quad 2 + 3 = z_3
\]
From the side condition, \(z_3 = 5\), so \(a^{\prime} = 5\). Hence,
\[
  \dfrac{2 + 3 \to 5}{1 + 2 + 3 \to 1 + 5}
\]
We then evaluate \(1 + 5\), which is similar to \(2 + 3\). By applying \textsc{PLUS}, we obtain:
\[
  1 + 5 \to 6
\]
\end{eg}

\chapter{Semantics Basics}

\section{Introduction to Semantics}
Consider the following C++14 program:
\begin{lstlisting}[language=iPython]
int main (int argc, char *argv[]) {
  int i = 1;
  int v = 2;
  i = i++ + v;
  printf("%d\n", i);

  return 0;
}
\end{lstlisting}

Using the \texttt{clang} compiler it prints 3, while using the \texttt{VS} compiler, it prints 4. However, what should it do and what is it supposed to do? The answer comes from the semantics of programming languages. 

Most programming languages have semantics described in natural language, which is preferable for many purposes but can be \textit{imprecise} and \textit{ambiguous}. For example, consider a programming language with a built-in function \texttt{sort}, where the documentation says:

\textit{The \texttt{sort} function takes an array and a comparator function as arguments. It sorts the array using the comparator function and returns it.}

This description is \textit{imprecise}, as there is no information on whether it modifies the original array or creates a new one, nor even the order of sorting. 

We can modify the description to be more precise and less ambiguous, yet natural language itself is hard to make fully precise. The ambiguity and imprecision of natural language can lead to diverging implementations of the same programming language, which leads to unexpected behavior.  

Therefore, we want formal semantics that give meaning to programs in a \textbf{precise}, \textbf{unambiguous}, and \textbf{concise} way. 

Different ways include:

- \textbf{Operational semantics}: how is a program executed on an abstract machine?

- \textbf{Denotational semantics}: what does the program mean in terms of mathematical objects?

- \textbf{Axiomatic semantics}: what should we be able to prove about this program?